\subsection{Uniform Continuity and Compactness}

\begin{definition} \label{def:uniform-continuity} \hlt{Uniform Continuity}\\
A mapping $f:X\to Y$ between metric spaces is \hlt{uniformly continuous} on $X$ if $\forall$ $\varepsilon>0$ $\exists$ $\delta>0$ such that $d_X(p,q)<\delta$ $\Longrightarrow$ $d_Y(f(p),f(q))<\varepsilon$ $\forall p,q\in X$. The same $\delta$ must work for every point of $X$.\\
\hlt{Structural Commentary}: Uniform continuity is setwise: for each $\varepsilon$, one $\delta$ must work for all pairs in the set.
\end{definition}

\begin{proposition} \label{prop:uniform-continuity-sequential-test} \hlt{Sequential Test for Uniform Continuity}\\
Let $X,Y$ be metric spaces and let $f:X\to Y$. The following are equivalent:
\begin{enumerate}[label=(\roman*)]
\setlength{\itemsep}{0pt}
\item $f$ is uniformly continuous on $X$;
\item whenever $(p_n)$ and $(q_n)$ are sequences in $X$ with $d_X(p_n,q_n)\to0$, one has $d_Y(f(p_n),f(q_n))\to0$.
\end{enumerate}
\begin{proof}
\begin{enumerate}[label=(\roman*).]
\setlength{\itemsep}{0pt}
\item[$(i)$]$\Rightarrow(ii)$ Given $\varepsilon>0$, choose $\delta>0$ such that $d_X(x,y)<\delta$ implies $d_Y(f(x),f(y))<\varepsilon$ for all $x,y\in X$.\\
Since $d_X(p_n,q_n)\to0$, eventually $d_X(p_n,q_n)<\delta$, hence eventually $d_Y(f(p_n),f(q_n))<\varepsilon$.
\item[$(ii)$]$\Rightarrow(i)$ Prove the contrapositive. If $f$ is not uniformly continuous, then $\exists \varepsilon_0>0$ such that $\forall n\geq1$, $\exists p_n,q_n\in X$ with $d_X(p_n,q_n)<1/n$ and $d_Y(f(p_n),f(q_n))\geq\varepsilon_0$.\\
Thus $d_X(p_n,q_n)\to0$ but $d_Y(f(p_n),f(q_n))\not\to0$, contradicting (ii).
\end{enumerate}
\hlt{Structural Commentary}: Uniform continuity preserves vanishing pairwise distances.\\
\hlt{Key Insight}: Failure of one global $\delta$ manufactures two close input sequences with separated outputs.
\end{proof}
\end{proposition}

\begin{figure}[H]
  \centering
  \begin{tikzpicture}[
  scale=0.95,
  >=Stealth,
  every node/.style={font=\scriptsize},
  axis/.style={->, gray!70, line width=0.45pt},
  curve/.style={blue!70!black, line width=0.9pt},
  guide/.style={gray!55, dashed, line width=0.4pt},
  epsmark/.style={orange!80!black, line width=0.65pt},
  bracket/.style={teal!75!black, line width=0.8pt},
  sample/.style={blue!70!black}
]
  \draw[axis] (-0.25,0) -- (4.00,0) node[right] {$x$};
  \draw[axis] (0,-0.45) -- (0,2.65) node[above] {$y$};

  \draw[curve, domain=0.66:3.45, samples=120] plot ({\x},{1/(\x*\x)});

  \foreach \x/\name/\d/\dlabel in {2.85/$x_1$/0.46/$\delta_1$,1.75/$x_2$/0.26/$\delta_2$,1.03/$x_3$/0.11/$\delta_3$} {
    \pgfmathsetmacro{\y}{1/(\x*\x)}
    \fill[sample] (\x,\y) circle (0.035);
    \draw[guide] (\x,0) -- (\x,\y);
    \node[above] at (\x,0.03) {\name};
    \draw[bracket] ({\x-\d},-0.12) -- ({\x+\d},-0.12);
    \draw[bracket] ({\x-\d},-0.18) -- ({\x-\d},-0.06);
    \draw[bracket] ({\x+\d},-0.18) -- ({\x+\d},-0.06);
    \node[teal!75!black, below] at (\x,-0.19) {\dlabel};
  }

  \draw[guide] (3.35,0.45) -- (3.70,0.45);
  \draw[guide] (3.35,1.20) -- (3.70,1.20);
  \draw[<->, epsmark] (3.60,0.45) -- (3.60,1.20);
  \node[orange!80!black, right] at (3.60,0.825) {$2\varepsilon$};
\end{tikzpicture}

  \caption{For $f(x)=1/x^2$ near $0$, a fixed vertical tolerance forces smaller pointwise horizontal tolerances.}
  \label{fig:uniform-continuity-delta-shrinks-one-over-x-squared}
\end{figure}

\begin{theorem} \label{thm:compact-continuity-uniform-continuity} \hlt{Continuity on Compact Sets Is Uniform}\\
If $f:X\to Y$ is continuous and $X$ is compact, then $f$ is uniformly continuous.
\begin{proof}
Let $\varepsilon>0$. For each $a\in X$, choose $r_a>0$ such that $d_X(x,a)<r_a$ implies $d_Y(f(x),f(a))<\varepsilon/2$.\\
Compactness gives a finite cover $X\subseteq\bigcup_{i=1}^n B_{r_{a_i}/2}(a_i)$. Set $\delta\defeq\frac12\min_i r_{a_i}>0$. If $d_X(p,q)<\delta$ and $p\in B_{r_{a_i}/2}(a_i)$, then $d_X(q,a_i)<r_{a_i}$ also; hence $f(p)$ and $f(q)$ are each within $\varepsilon/2$ of $f(a_i)$. Thus $d_Y(f(p),f(q))<\varepsilon$.\\
\hlt{Structural Commentary}: Compactness converts pointwise radii into finitely many radii with positive minimum.\\
\hlt{Key Insight}: Uniform continuity is continuity with the point-dependence removed by a finite subcover.
\end{proof}
\end{theorem}

\begin{theorem} \label{thm:compactness-hypotheses-are-essential} \hlt{Compactness Hypotheses Are Essential}\\
Let $E$ be a noncompact subset of $\R$. Then:
\begin{enumerate}[label=(\roman*)]
\setlength{\itemsep}{0pt}
\item there exists a continuous function on $E$ which is not bounded;
\item there exists a continuous bounded function on $E$ which has no maximum;
\item if $E$ is bounded, then there exists a continuous function on $E$ which is not uniformly continuous.
\end{enumerate}
\begin{proof}
\begin{enumerate}[label=(\roman*)]
\setlength{\itemsep}{0pt}
\item If $E$ is unbounded, $f(x)=x$ works. If $E$ is bounded, Theorem \ref{thm:heine-borel-rk} gives $x_0\in\overline{E}\setminus E$, and $f(x)=1/(x-x_0)$ is continuous on $E$ and unbounded.
\item If $E$ is unbounded, $h(x)=x^2/(1+x^2)$ has $\sup_E h=1$ but $h<1$. If $E$ is bounded, with $x_0$ as above, $g(x)=1/(1+(x-x_0)^2)$ is bounded by $1$ and has supremum $1$ but no maximum.
\item If $E$ is bounded, let $f(x)=1/(x-x_0)$. Given $\delta>0$, let $x\in E$ with $\abs{x-x_0}<\delta/2$, then $y\in E$ with $\abs{y-x_0}<\min\{\delta/2,1/(\abs{f(x)}+1)\}$. Then $\abs{x-y}<\delta$, $\abs{f(x)-f(y)}>1$, so $f$ is not uniformly continuous.
\end{enumerate}
\hlt{Structural Commentary}: Compactness is essential in the preceding theorems.\\
\hlt{Key Insight}: Noncompact sets either escape to infinity or miss an accumulation point.
\end{proof}
\end{theorem}
